\subsubsection{\texttt{dotnet-coverage}}\label{dotnet-coverage}

Pokrivenost koda meri koji delovi programa su izvršeni tokom pokretanja testova. To nije dokaz ispravnosti programa. Korisna je kao povratna informacija: pokazuje grane koje testovi ne dodiruju i pomaže da se razgovor o nedostajućim testovima veže za konkretan kod.

Alat \texttt{dotnet-coverage} sakuplja pokrivenost za .NET proces koji pokreće zadatu komandu. U ovoj vežbi ga koristimo kao lokalni .NET alat i izvozimo rezultat u Cobertura XML format. Cobertura je pogodan za pregled u alatima za izveštavanje i CI sistemima, ali u ovoj vežbi ne postavljamo prag koji ruši proveru.

Manifest alata je deo direktorijuma vežbe:

\codefile{.config/dotnet-tools.json}{json}

Primer koristi malu biblioteku sa granama u logici popusta. Cilj nije domen prodaje, već kratak kod nad kojim se jasno vidi razlika između prolaznih testova i dobrih testova.

\codefile{CoverageLab/CoverageLab.csproj}{xml}

\codefile{CoverageLab/CustomerKind.cs}{csharp}

\codefile{CoverageLab/DiscountPolicy.cs}{csharp}

Prvi skup testova proverava deo ponašanja. Svi testovi prolaze, ali pokrivenost pokazuje da nisu provereni partnerski popust i obrada neispravnog međuzbira.

\codefile{CoverageLab.Initial.Tests/CoverageLab.Initial.Tests.csproj}{xml}

\codefile{CoverageLab.Initial.Tests/DiscountPolicyTests.cs}{csharp}

Drugi skup testova dodaje propuštene slučajeve. On ne menja produkcioni kod, već pokazuje kako ista biblioteka dobija bolju pokrivenost kada testovi eksplicitno pokriju odluke u metodi.

\codefile{CoverageLab.Complete.Tests/CoverageLab.Complete.Tests.csproj}{xml}

\codefile{CoverageLab.Complete.Tests/DiscountPolicyTests.cs}{csharp}

Vežba se pokreće iz direktorijuma \texttt{04\_dotnet\_coverage}:

\begin{minted}{bash}
$ dotnet tool restore
$ dotnet restore CoverageLab.sln
$ dotnet test --solution CoverageLab.sln --no-restore
\end{minted}

Komande za pokrivenost ispod koriste statičku instrumentaciju preko opcije \texttt{-if}. Instrumentiše se biblioteka koja je kopirana u izlazni direktorijum konkretnog test projekta, pa pre ovih komandi treba jednom pokrenuti \texttt{dotnet test}.

Pokrivenost početnog skupa testova:

\begin{minted}{bash}
$ dotnet tool run dotnet-coverage -- collect \
    "dotnet CoverageLab.Initial.Tests/bin/Debug/net10.0/CoverageLab.Initial.Tests.dll" \
    -if CoverageLab.Initial.Tests/bin/Debug/net10.0/CoverageLab.dll \
    -f cobertura \
    -o coverage-initial.cobertura.xml
\end{minted}

Pokrivenost dopunjenog skupa testova:

\begin{minted}{bash}
$ dotnet tool run dotnet-coverage -- collect \
    "dotnet CoverageLab.Complete.Tests/bin/Debug/net10.0/CoverageLab.Complete.Tests.dll" \
    -if CoverageLab.Complete.Tests/bin/Debug/net10.0/CoverageLab.dll \
    -f cobertura \
    -o coverage-complete.cobertura.xml
\end{minted}

U dobijenim XML datotekama treba uporediti vrednosti atributa \texttt{line-rate} i \texttt{branch-rate}. Bolji rezultat nije cilj sam po sebi. Ako test samo izvrši liniju bez smislenog tvrđenja, pokrivenost raste, ali kvalitet provere ne raste. Zato izveštaj treba čitati zajedno sa test kodom.

\texttt{Coverlet} je česta alternativa kada je potrebno da sakupljanje pokrivenosti bude direktno vezano za \texttt{dotnet test} i test adaptere. U ovoj vežbi koristimo \texttt{dotnet-coverage}, jer isti tok rada može da obuhvati bilo koju .NET komandu ili proces koji se pokreće iz komandne linije.
